\documentclass[../../main.tex]{subfiles}
\graphicspath{{\subfix{../../image/}}}

\begin{document}

\section{极值原理}\label{section:ch06-极值原理}

本节发展二阶椭圆型偏微分方程的极值原理。

极值原理方法基于如下观察：若 $C^{2}$ 函数 $u$ 在开集 $U$ 上于某点 $x_{0}\in U$ 达到最大值，则
\begin{align}
Du(x_{0})=0,\qquad D^{2}u(x_{0})\leqslant 0,
\end{align}
后一不等式意味着对称矩阵 $D^{2}u=((u_{x_{i}x_{j}}))$ 在 $x_{0}$ 处非正定。基于此的推论因而具有逐点特征，与\cref{subsection:ch06-弱解}--\cref{section:ch06-正则性} 中阐述的基于积分的能量方法完全不同。

此外我们将需要要求解 $u$ 至少为 $C^{2}$，以使考虑 $Du$、$D^{2}u$ 的逐点值有意义。（然而鉴于\cref{section:ch06-正则性} 的正则性理论，我们知道弱解至少在系数充分正则时确实如此光滑。）我们很快将看到，现在考虑具有非散度型的椭圆算子 $L$ 也最为合适：
\begin{align}
Lu=-\sum_{i,j=1}^{n}a_{ij}u_{x_{i}x_{j}}+\sum_{i=1}^{n}b_{i}u_{x_{i}}+cu,
\end{align}
其中系数 $a_{ij}$、$b_{i}$、$c$ 连续，且——如通常那样——\cref{definition:ch06-椭圆性} 的一致椭圆性条件成立。我们还不失一般性地继续假设对称性条件 $a_{ij}=a_{ji}$（$i,j=1,\cdots,n$）。

\subsection{弱极值原理}\label{subsection:ch06-弱极值}

我们首先确定函数必须在边界上达到最大（或最小）值的情形。我们总假设 $U\subset\mathbb{R}^{n}$ 开、有界。

\begin{theorem}[弱极值原理]\label{theorem:ch06-弱极值}设 $u\in C^{2}(U)\cap C(\overline{U})$，且 $c\equiv 0$ 于 $U$。
\begin{enumerate}[(i)]
\item 若
\begin{align}
Lu\leqslant 0\quad\text{于}\ U,
\end{align}
则
\begin{align*}
\max_{\overline{U}}u=\max_{\partial U}u.
\end{align*}
\item 若
\begin{align}
Lu\geqslant 0\quad\text{于}\ U,
\end{align}
则
\begin{align*}
\min_{\overline{U}}u=\min_{\partial U}u.
\end{align*}
\end{enumerate}
\end{theorem}

\begin{remark}\label{remark:ch06-上下解}
满足式~\eqref{align:ch06-下解} 的函数称为下解。我们因而断言下解在 $\partial U$ 上达到其最大值。类似地，若式~\eqref{align:ch06-上解} 成立，$u$ 为上解，在 $\partial U$ 上达到其最小值。
\end{remark}

\begin{proof}
1. 先假设我们有严格不等式
\begin{align}
Lu<0\quad\text{于}\ U,
\end{align}
而存在点 $x_{0}\in U$ 使
\begin{align}
u(x_{0})=\max_{\overline{U}}u.
\end{align}
现在此最大点 $x_{0}$ 处，我们有
\begin{align}
Du(x_{0})=0
\end{align}
且
\begin{align}
D^{2}u(x_{0})\leqslant 0.
\end{align}

2. 由于矩阵 $A=((a_{ij}(x_{0})))$ 对称正定，存在正交矩阵 $O=((o_{ij}))$ 使
\begin{align}
OAO^{T}=\operatorname{diag}(d_{1},\cdots,d_{n}),\qquad OO^{T}=I,
\end{align}
其中 $d_{k}>0$（$k=1,\cdots,n$）。写 $y=x_{0}+O(x-x_{0})$。则 $x-x_{0}=O^{T}(y-x_{0})$，故
\begin{align*}
u_{x_{i}}=\sum_{k=1}^{n}u_{y_{k}}o_{ki},\qquad
u_{x_{i}x_{j}}=\sum_{k,l=1}^{n}u_{y_{k}y_{l}}o_{ki}o_{lj}\qquad (i,j=1,\cdots,n).
\end{align*}
因而在点 $x_{0}$ 处，
\begin{align}
\sum_{i,j=1}^{n}a_{ij}u_{x_{i}x_{j}}
=\sum_{k,l=1}^{n}\sum_{i,j=1}^{n}a_{ij}u_{y_{k}y_{l}}o_{ki}o_{lj}
=\sum_{k=1}^{n}d_{k}u_{y_{k}y_{k}}\qquad\text{由上式}
\leqslant 0,
\end{align}
因为 $d_{k}>0$ 且由式~\eqref{align:ch06-正交} 有 $u_{y_{k}y_{k}}(x_{0})\leqslant 0$（$k=1,\cdots,n$）。

3. 于是在 $x_{0}$ 处，由前述两式，
\begin{align*}
Lu=-\sum_{i,j=1}^{n}a_{ij}u_{x_{i}x_{j}}+\sum_{i=1}^{n}b_{i}u_{x_{i}}\geqslant 0.
\end{align*}
故式~\eqref{align:ch06-严格小于} 与~\eqref{align:ch06-内最大} 不相容，矛盾。

4. 在式~\eqref{align:ch06-下解} 成立的一般情形，写
\begin{align*}
u_{\varepsilon}(x)\triangleq u(x)+\varepsilon e^{\lambda x_{1}}\qquad (x\in U),
\end{align*}
其中 $\lambda>0$ 将在下面选取，$\varepsilon>0$。回忆（如\cref{theorem:ch06-边界H2正则} 的证明中那样）一致椭圆性条件蕴含 $a_{ii}(x)\geqslant\theta$（$i=1,\cdots,n$，$x\in U$）。因此
\begin{align*}
Lu_{\varepsilon}=Lu+\varepsilon L(e^{\lambda x_{1}})
\leqslant\varepsilon e^{\lambda x_{1}}[-\lambda^{2}a_{11}+\lambda b_{1}]
\leqslant\varepsilon e^{\lambda x_{1}}[-\lambda^{2}\theta+\|b\|_{L^{\infty}}\lambda]
<0\quad\text{于}\ U,
\end{align*}
只要选 $\lambda>0$ 充分大。然后由上述第 1、2 步，$\max_{\overline{U}}u_{\varepsilon}=\max_{\partial U}u_{\varepsilon}$。令 $\varepsilon\to 0$，得 $\max_{\overline{U}}u=\max_{\partial U}u$。这就证明了 (i)。

5. 由于每当 $u$ 为上解时 $-u$ 为下解，断言 (ii) 随之成立。
\end{proof}

下面我们修改极值原理以允许非零零阶系数 $c$。回忆\cref{section:appA-函数记号}，$u^{+}=\max(u,0)$、$u^{-}=-\min(u,0)$。

\begin{theorem}[$c\geqslant 0$ 的弱极值原理]\label{theorem:ch06-c非负弱极值}设 $u\in C^{2}(U)\cap C(\overline{U})$，且
\begin{align*}
c\geqslant 0\quad\text{于}\ U.
\end{align*}
\begin{enumerate}[(i)]
\item 若 $Lu\leqslant 0$ 于 $U$，则
\begin{align}
\max_{\overline{U}}u\leqslant\max_{\partial U}u^{+}.
\end{align}
\item 类似地，若 $Lu\geqslant 0$ 于 $U$，则
\begin{align}
\min_{\overline{U}}u\geqslant-\max_{\partial U}u^{-}.
\end{align}
\end{enumerate}
\end{theorem}

\begin{remark}\label{remark:ch06-c非负弱极值注}
特别地，若 $Lu=0$ 于 $U$，则
\begin{align}
\max_{\overline{U}}|u|=\max_{\partial U}|u|.
\end{align}
\end{remark}

\begin{proof}
1. 设 $u$ 为下解，令 $V\triangleq\{x\in U\mid u(x)>0\}$。则
\begin{align*}
Ku\triangleq Lu-cu\leqslant-cu\leqslant 0\quad\text{于}\ V.
\end{align*}
算子 $K$ 没有零阶项，因而\cref{theorem:ch06-弱极值} 蕴含 $\max_{\overline{V}}u=\max_{\partial V}u=\max_{\partial U}u^{+}$。这就在 $V\neq\emptyset$ 的情形给出式~\eqref{align:ch06-上解}。否则 $u\leqslant 0$ 处处于 $U$，式~\eqref{align:ch06-上解} 同样成立。

2. 一旦我们注意 $(-u)^{+}=u^{-}$，断言 (ii) 由 (i) 应用于 $-u$ 即得。
\end{proof}

\subsection{强极值原理}\label{subsection:ch06-强极值}

下面我们通过证明下解 $u$ 根本不能在连通区域的内点处达到其最大值（除非 $u$ 为常数），来实质性地加强前述断言。这一陈述即强极值原理，它依赖于对边界最大点处外法向导数 $\dfrac{\partial u}{\partial\nu}$ 的如下精细分析。

\begin{lemma}[Hopf 引理]\label{lemma:ch06-Hopf}设 $u\in C^{2}(U)\cap C^{1}(\overline{U})$，且 $c\equiv 0$ 于 $U$。另设
\begin{align*}
Lu\leqslant 0\quad\text{于}\ U,
\end{align*}
且存在点 $x_{0}\in\partial U$ 使
\begin{align}
u(x_{0})>u(x)\quad\text{对一切}\ x\in U.
\end{align}
最后设 $U$ 在 $x_{0}$ 处满足内球条件；即存在开球 $B\subset U$ 使 $x_{0}\in\partial B$。
\begin{enumerate}[(i)]
\item 则
\begin{align*}
\frac{\partial u}{\partial\nu}(x_{0})>0,
\end{align*}
其中 $\nu$ 为 $B$ 在 $x_{0}$ 处的外单位法向。
\item 若 $c\geqslant 0$ 于 $U$，则只要 $u(x_{0})\geqslant 0$，同一结论成立。
\end{enumerate}
\end{lemma}

(i) 的重要性在于严格不等式：$\dfrac{\partial u}{\partial\nu}(x_{0})\geqslant 0$ 是显然的。注意若 $\partial U$ 为 $C^{2}$，内球条件自动成立。

\begin{proof}
1. 设 $c\geqslant 0$。我们不妨进一步设 $B=B^{0}(0,r)$（某半径 $r>0$）。定义
\begin{align}
v(x)\triangleq e^{-\lambda|x|^{2}}-e^{-\lambda r^{2}}\qquad (x\in B(0,r))
\end{align}
其中 $\lambda>0$ 如下选取。然后利用一致椭圆性条件，计算
\begin{align*}
Lv&=-\sum_{i,j=1}^{n}a_{ij}v_{x_{i}x_{j}}+\sum_{i=1}^{n}b_{i}v_{x_{i}}+cv\\
&=e^{-\lambda|x|^{2}}\sum_{i,j=1}^{n}a_{ij}(-4\lambda^{2}x_{i}x_{j}+2\lambda\delta_{ij})
-e^{-\lambda|x|^{2}}\sum_{i=1}^{n}b_{i}2\lambda x_{i}
+c(e^{-\lambda|x|^{2}}-e^{-\lambda r^{2}})\\
&\leqslant e^{-\lambda|x|^{2}}(-4\theta\lambda^{2}|x|^{2}+2\lambda\operatorname{tr}A+2\lambda|b||x|+c),
\end{align*}
其中 $A=((a_{ij}))$、$b=(b_{1},\cdots,b_{n})$。其次考虑开环形区域 $R\triangleq B^{0}(0,r)\setminus B(0,r/2)$。我们有
\begin{align}
Lv\leqslant e^{-\lambda|x|^{2}}(-\theta\lambda^{2}r^{2}+2\lambda\operatorname{tr}A+2\lambda|b|r+c)\leqslant 0\quad\text{于}\ R,
\end{align}
只要 $\lambda>0$ 固定充分大。

2. 由式~\eqref{align:ch06-Hopf条件}，存在常数 $\varepsilon>0$ 充分小，使
\begin{align}
u(x_{0})\geqslant u(x)+\varepsilon v(x)\qquad (x\in\partial B(0,r/2)).
\end{align}
此外注意
\begin{align}
u(x_{0})\geqslant u(x)+\varepsilon v(x)\qquad (x\in\partial B(0,r)),
\end{align}
因为 $v\equiv 0$ 于 $\partial B(0,r)$。

3. 由式~\eqref{align:ch06-HopfLv}，我们看到
\begin{align*}
L(u+\varepsilon v-u(x_{0}))\leqslant-cu(x_{0})\leqslant 0\quad\text{于}\ R,
\end{align*}
且由前述两式，我们观察
\begin{align*}
u+\varepsilon v-u(x_{0})\leqslant 0\quad\text{于}\ \partial R.
\end{align*}
由弱极值原理\cref{theorem:ch06-c非负弱极值}，$u+\varepsilon v-u(x_{0})\leqslant 0$ 于 $R$。但 $u(x_{0})+\varepsilon v(x_{0})-u(x_{0})=0$，故
\begin{align*}
\frac{\partial u}{\partial\nu}(x_{0})+\varepsilon\frac{\partial v}{\partial\nu}(x_{0})\geqslant 0.
\end{align*}
因而
\begin{align*}
\frac{\partial u}{\partial\nu}(x_{0})\geqslant-\varepsilon\frac{\partial v}{\partial\nu}(x_{0})
=-\frac{\varepsilon}{r}Dv(x_{0})\cdot x_{0}
=2\lambda\varepsilon r e^{-\lambda r^{2}}>0,
\end{align*}
如所要求。
\end{proof}

Hopf 引理是下面证明中的主要技术工具：

\begin{theorem}[强极值原理]\label{theorem:ch06-强极值}设 $u\in C^{2}(U)\cap C(\overline{U})$，且 $c\equiv 0$ 于 $U$。另设 $U$ 连通、开、有界。
\begin{enumerate}[(i)]
\item 若 $Lu\leqslant 0$ 于 $U$，且 $u$ 在 $\overline{U}$ 上于一内点达到其最大值，则 $u$ 在 $U$ 内为常数。
\item 类似地，若 $Lu\geqslant 0$ 于 $U$，且 $u$ 在 $\overline{U}$ 上于一内点达到其最小值，则 $u$ 在 $U$ 内为常数。
\end{enumerate}
\end{theorem}

\begin{proof}
写 $M\triangleq\max_{\overline{U}}u$，$C\triangleq\{x\in U\mid u(x)=M\}$。则若 $u\not\equiv M$，令 $V\triangleq\{x\in U\mid u(x)<M\}$。选一点 $y\in V$ 满足 $\operatorname{dist}(y,C)<\operatorname{dist}(y,\partial U)$，并令 $B$ 表示以 $y$ 为中心、其内部位于 $V$ 中的最大球。于是存在某点 $x_{0}\in C$，且 $x_{0}\in\partial B$。显然 $V$ 在 $x_{0}$ 处满足内球条件，因而 Hopf 引理 (i) 蕴含 $\dfrac{\partial u}{\partial\nu}(x_{0})>0$。但这是矛盾：因为 $u$ 在 $x_{0}\in U$ 处达到最大值，我们有 $Du(x_{0})=0$。
\end{proof}

若零阶项 $c$ 非负，我们有强极值原理的这一版本：

\begin{theorem}[$c\geqslant 0$ 的强极值原理]\label{theorem:ch06-c非负强极值}设 $u\in C^{2}(U)\cap C(\overline{U})$，且 $c\geqslant 0$ 于 $U$。另设 $U$ 连通。
\begin{enumerate}[(i)]
\item 若 $Lu\leqslant 0$ 于 $U$，且 $u$ 在 $\overline{U}$ 上于一内点达到一非负最大值，则 $u$ 在 $U$ 内为常数。
\item 类似地，若 $Lu\geqslant 0$ 于 $U$，且 $u$ 在 $\overline{U}$ 上于一内点达到一非正最小值，则 $u$ 在 $U$ 内为常数。
\end{enumerate}
\end{theorem}

证明如上，只是我们使用 Hopf 引理的 (ii)。

\subsection{Harnack 不等式}\label{subsection:ch06-Harnack}

Harnack 不等式断言：非负解的值是可比的，至少在任一远离边界的子区域中如此。我们如通常那样假设
\begin{align*}
Lu=-\sum_{i,j=1}^{n}a_{ij}u_{x_{i}x_{j}}+\sum_{i=1}^{n}b_{i}u_{x_{i}}+cu.
\end{align*}

\begin{theorem}[Harnack 不等式]\label{theorem:ch06-Harnack}设 $u\geqslant 0$ 为 $Lu=0$ 于 $U$ 中的 $C^{2}$ 解，且设 $V\subset \subset U$ 连通。则存在常数 $C$，使
\begin{align}
\sup_{V}u\leqslant C\inf_{V}u.
\end{align}
常数 $C$ 只依赖于 $V$ 及 $L$ 的系数。
\end{theorem}

若系数仅为有界可测，此断言仍成立：见 Gilbarg--Trudinger [G-T]。然而我们将仅对 $b_{i}\equiv c\equiv 0$ 且 $a_{ij}$ 光滑（$i,j=1,\cdots,n$）的特殊情形提供证明。

\begin{proof}
1. 我们可设 $u>0$ 于 $U$，否则可把结果应用于 $u+\varepsilon$ 然后令 $\varepsilon\to 0^{+}$。令
\begin{align}
v\triangleq\log u.
\end{align}
由于 $Lu=0$，我们计算
\begin{align}
\sum_{i,j=1}^{n}a_{ij}v_{x_{i}x_{j}}+a_{ij}v_{x_{i}}v_{x_{j}}=0\quad\text{于}\ U.
\end{align}
定义
\begin{align}
w\triangleq\sum_{i,j=1}^{n}a_{ij}v_{x_{i}}v_{x_{j}},
\end{align}
于是上式说
\begin{align}
-\sum_{i,j=1}^{n}a_{ij}v_{x_{i}x_{j}}=w.
\end{align}

2. 我们对 $k,l=1,\cdots,n$ 计算
\begin{align*}
w_{x_{k}x_{l}}=\sum_{i,j=1}^{n}2a_{ij}v_{x_{i}x_{k}x_{l}}v_{x_{j}}+2a_{ij}v_{x_{i}x_{k}}v_{x_{j}x_{l}}+R,
\end{align*}
其中余项 $R$（由落在系数上的导数产生）满足形如
\begin{align}
|R|\leqslant\varepsilon|D^{2}v|^{2}+C(\varepsilon)|Dv|^{2}
\end{align}
的估计，对每个 $\varepsilon>0$。因而
\begin{align}
-\sum_{k,l=1}^{n}a_{kl}w_{x_{k}x_{l}}
=2\sum_{i,j=1}^{n}a_{ij}v_{x_{j}}\left(-\sum_{k,l=1}^{n}a_{lk}v_{x_{i}x_{k}x_{l}}\right)
-2\sum_{i,j,k,l=1}^{n}a_{ij}a_{kl}v_{x_{i}x_{k}}v_{x_{j}x_{l}}-R.
\end{align}
对上式微分，我们看到
\begin{align}
-\sum_{k,l=1}^{n}a_{lk}v_{x_{i}x_{k}x_{l}}=w_{x_{i}}+R_{i}\qquad (i=1,\cdots,n),
\end{align}
其中 $R_{i}$ 表示另一余项，满足 $|R_{i}|\leqslant C|D^{2}v|$。此外，一致椭圆性条件蕴含
\begin{align}
\sum_{i,j,k,l=1}^{n}a_{ij}a_{kl}v_{x_{i}x_{k}}v_{x_{j}x_{l}}\geqslant\theta^{2}|D^{2}v|^{2}.
\end{align}
把前述两式代入上式，然后选 $\varepsilon>0$ 充分小，我们导出微分不等式
\begin{align}\label{align:ch06-Harnack微分不等式}
-\sum_{k,l=1}^{n}a_{kl}w_{x_{k}x_{l}}+\sum_{k=1}^{n}\widetilde{b}_{k}w_{x_{k}}\leqslant-\frac{\theta^{2}}{2}|D^{2}v|^{2}+C|Dv|^{2},
\end{align}
其中
\begin{align}
\widetilde{b}_{k}\triangleq-2\sum_{l=1}^{n}a_{kl}v_{x_{l}}\qquad (k=1,\cdots,n).
\end{align}

3. 现设 $V\subset \subset U$ 为半径 $r>0$ 的开球。选截断函数 $\zeta\in C^{\infty}(U)$ 使 $0\leqslant\zeta\leqslant 1$、$\zeta=0$ 于 $\partial U$、$\zeta\equiv 1$ 于 $V$。令
\begin{align}
z\triangleq\zeta^{4}w.
\end{align}
我们现假设 $z$ 在某点 $x_{0}\in U$ 达到其最大值。于是在 $x_{0}$ 处有
\begin{align}
\zeta w_{x_{k}}+4\zeta_{x_{k}}w=0\qquad (k=1,\cdots,n).
\end{align}
此外，在此点处我们有不等式
\begin{align*}
0\leqslant-\sum_{k,l=1}^{n}a_{kl}z_{x_{k}x_{l}}+\sum_{k=1}^{n}\widetilde{b}_{k}z_{x_{k}}
=-\sum_{k,l=1}^{n}a_{kl}(\zeta^{4}w)_{x_{k}x_{l}}+\sum_{k=1}^{n}\widetilde{b}_{k}(\zeta^{4}w)_{x_{k}}.
\end{align*}
因而
\begin{align}
0\leqslant\zeta^{4}\left(-\sum_{k,l=1}^{n}a_{kl}w_{x_{k}x_{l}}+\sum_{k=1}^{n}\widetilde{b}_{k}w_{x_{k}}\right)+\widehat{R},
\end{align}
其中余项 $\widehat{R}$（由落在截断函数 $\zeta$ 上的导数组成）满足估计
\begin{align*}
|\widehat{R}|\leqslant C(\zeta^{2}w+\zeta^{3}|Dw|).
\end{align*}
现回忆式~\eqref{align:ch06-Harnack最大条件}, 我们看到事实上
\begin{align}
|\widehat{R}|\leqslant C\zeta^{2}w.
\end{align}
把此估计与式~\eqref{align:ch06-Harnack微分不等式} 代入上式：
\begin{align*}
\zeta^{4}|D^{2}v|^{2}\leqslant C\zeta^{4}|Dv|^{2}+C\zeta^{2}w.
\end{align*}
但 $\theta|Dv|^{2}\leqslant w$，且由式~\eqref{align:ch06-w定义} 还有 $w\leqslant C|D^{2}v|$。由此得
\begin{align*}
\zeta^{4}w^{2}\leqslant C\zeta^{2}w,
\end{align*}
因而在点 $x_{0}$ 处
\begin{align*}
z=\zeta^{4}w\leqslant C.
\end{align*}
由于 $z$ 在 $x_{0}$ 处达到最大值且 $\zeta\equiv 1$ 于 $V$，我们因而得到估计
\begin{align}
|Dv|\leqslant C\quad\text{于区域}\ V\ \text{内}.
\end{align}
于是在 $V$ 内沿连接任意两点的路径积分，得 $\sup_{V}u\leqslant C\inf_{V}u$。一般连通 $V\subset \subset U$ 的情形由有限个球覆盖即得。
\end{proof}

\end{document}
